\documentclass[10pt,a4paper]{amsart}
\usepackage{amsmath,amssymb,amsthm,graphicx}
\usepackage{verbatim}
\newenvironment{reference}{\comment}{\endcomment}
%\newenvironment{reference}{}{}
\newenvironment{slogan}{\comment}{\endcomment}
\newenvironment{history}{\comment}{\endcomment}
% For commutative diagrams we use Xy-pic
\usepackage[all]{xy}

% We use 2cell for 2-commutative diagrams.
\xyoption{2cell}
% Theorem environments.
%
\theoremstyle{plain}
\newtheorem{theorem}[subsection]{Theorem}
\newtheorem{proposition}[subsection]{Proposition}
\newtheorem{lemma}[subsection]{Lemma}

\theoremstyle{definition}
\newtheorem{definition}[subsection]{Definition}
\newtheorem{example}[subsection]{Example}
\newtheorem{exercise}[subsection]{Exercise}
\newtheorem{situation}[subsection]{Situation}

\theoremstyle{remark}
\newtheorem{remark}[subsection]{Remark}
\newtheorem{remarks}[subsection]{Remarks}

\numberwithin{equation}{subsection}

% Macros
%
\def\lim{\mathop{\mathrm{lim}}\nolimits}
\def\colim{\mathop{\mathrm{colim}}\nolimits}
\def\Spec{\mathop{\mathrm{Spec}}}
\def\Hom{\mathop{\mathrm{Hom}}\nolimits}
\def\Ext{\mathop{\mathrm{Ext}}\nolimits}
\def\SheafHom{\mathop{\mathcal{H}\!\mathit{om}}\nolimits}
\def\SheafExt{\mathop{\mathcal{E}\!\mathit{xt}}\nolimits}
\def\Sch{\mathit{Sch}}
\def\Mor{\mathop{\mathrm{Mor}}\nolimits}
\def\Ob{\mathop{\mathrm{Ob}}\nolimits}
\def\Sh{\mathop{\mathit{Sh}}\nolimits}
\def\NL{\mathop{N\!L}\nolimits}
\def\CH{\mathop{\mathrm{CH}}\nolimits}
\def\proetale{{pro\text{-}\acute{e}tale}}
\def\etale{{\acute{e}tale}}
\def\QCoh{\mathit{QCoh}}
\def\Ker{\mathop{\mathrm{Ker}}}
\def\Im{\mathop{\mathrm{Im}}}
\def\Coker{\mathop{\mathrm{Coker}}}
\def\Coim{\mathop{\mathrm{Coim}}}

% Boxtimes
%
\DeclareMathSymbol{\boxtimes}{\mathbin}{AMSa}{"02}

%
% Macros for moduli stacks/spaces
%
\def\QCohstack{\mathcal{QC}\!\mathit{oh}}
\def\Cohstack{\mathcal{C}\!\mathit{oh}}
\def\Spacesstack{\mathcal{S}\!\mathit{paces}}
\def\Quotfunctor{\mathrm{Quot}}
\def\Hilbfunctor{\mathrm{Hilb}}
\def\Curvesstack{\mathcal{C}\!\mathit{urves}}
\def\Polarizedstack{\mathcal{P}\!\mathit{olarized}}
\def\Complexesstack{\mathcal{C}\!\mathit{omplexes}}
% \Pic is the operator that assigns to X its picard group, usage \Pic(X)
% \Picardstack_{X/B} denotes the Picard stack of X over B
% \Picardfunctor_{X/B} denotes the Picard functor of X over B
\def\Pic{\mathop{\mathrm{Pic}}\nolimits}
\def\Picardstack{\mathcal{P}\!\mathit{ic}}
\def\Picardfunctor{\mathrm{Pic}}
\def\Deformationcategory{\mathcal{D}\!\mathit{ef}}
\usepackage{cleveref}
\crefname{originallemma}{original lemma}{original lemmas}
\Crefname{originallemma}{Original lemma}{Original lemmas}
\makeatletter
\newlabel{lemma-base-change-quasi-compact}{{lemma-base-change-quasi-compact}{source}{}{sourcepointer-1}{}}
\newlabel{lemma-base-change-quasi-compact@cref}{{[originallemma][0][]lemma-base-change-quasi-compact}{[1][0][]source}{}{}{}}
\makeatother
\crefname{originallemma}{original lemma}{original lemmas}
\Crefname{originallemma}{Original lemma}{Original lemmas}
\makeatletter
\newlabel{lemma-gerbe-residual-gerbe-exists}{{lemma-gerbe-residual-gerbe-exists}{source}{}{sourcepointer-2}{}}
\newlabel{lemma-gerbe-residual-gerbe-exists@cref}{{[originallemma][0][]lemma-gerbe-residual-gerbe-exists}{[1][0][]source}{}{}{}}
\makeatother
\crefname{originallemma}{original lemma}{original lemmas}
\Crefname{originallemma}{Original lemma}{Original lemmas}
\makeatletter
\newlabel{lemma-monomorphism-cartesian-square-inertia}{{lemma-monomorphism-cartesian-square-inertia}{source}{}{sourcepointer-3}{}}
\newlabel{lemma-monomorphism-cartesian-square-inertia@cref}{{[originallemma][0][]lemma-monomorphism-cartesian-square-inertia}{[1][0][]source}{}{}{}}
\makeatother
\crefname{originalproposition}{original proposition}{original propositions}
\Crefname{originalproposition}{Original proposition}{Original propositions}
\makeatletter
\newlabel{proposition-open-stratum}{{proposition-open-stratum}{source}{}{sourcepointer-4}{}}
\newlabel{proposition-open-stratum@cref}{{[originalproposition][0][]proposition-open-stratum}{[1][0][]source}{}{}{}}
\makeatother
\crefname{originalreference}{original reference}{original references}
\Crefname{originalreference}{Original reference}{Original references}
\makeatletter
\newlabel{stacks-properties-lemma-reduced-closed-substack}{{stacks-properties-lemma-reduced-closed-substack}{source}{}{sourcepointer-5}{}}
\newlabel{stacks-properties-lemma-reduced-closed-substack@cref}{{[originalreference][0][]stacks-properties-lemma-reduced-closed-substack}{[1][0][]source}{}{}{}}
\makeatother
\title{Every point of a quasi-separated stack has a residual gerbe}
\author{The Stacks Project Authors}
\date{stacks-project:stacks-morphisms.tex:local-snapshot-20261012; GNU FDL 1.2 or later}
\begin{document}
\maketitle
\noindent\textit{Editorial note.} All delivered mathematical text is the author's own, copied from the pinned Stacks Project chapter source sources/stacks-algebraic-stacks-morphisms/pinned-original-source.tex; no statement or proof step was rewritten, reordered or removed.\par\medskip
\noindent\textit{Editorial note.} Preamble provenance: the Stacks Project keeps one single shared preamble for every chapter; this card takes it directly from the staged sources/stacks-algebraic-stacks-morphisms/preamble.tex as three verbatim spans (lines 8-12, 14-18 and 154-217), so no derived file was written and nothing was added or redefined. The wrapper class and the base packages are supplied by the assembler.\par\medskip
\noindent\textit{Editorial note.} Bibliography provenance: neither the counted statement nor its proof contains a citation (this lane rescans the whole span before writing the card), so the card carries no bibliography and no entry was reconstructed.\par\medskip
\noindent\textit{Editorial note.} Selection provenance: this card was opened under the file-based control contract (CONTROL.md), with the mechanical pre-card gate run on its statement and proof interval before the card was written.\par\medskip
\noindent\textit{Allowed standard prerequisites (editorial).}\begin{itemize}
\item The construction of the residual gerbe at a point of finite type through the closure of the point. Used inside the counted proof (source lines 7059-7070). Produces the residual gerbe for a point of finite type after localising.
\item The existence of points of finite type in every nonempty locally closed subset. Invoked inside the counted proof (source lines 7071-7076). Ensures the construction applies to the localisation at the given point.
\end{itemize}\medskip
\ifdefined\hypertarget\hypertarget{sourcepointer-1}{}\fi
\noindent\textit{Source reference.} Original lemma lemma-base-change-quasi-compact: sources/stacks-algebraic-stacks-morphisms/pinned-original-source.tex, source lines 1984--1984. Navigation pointer only: the counted text refers to the Stacks item lemma-base-change-quasi-compact, which is an existing labeled result of this same chapter and is not part of this card. The pointer carries the target's own label as its number and its definition line as its span; its content is not reproduced here.\par\medskip
\ifdefined\hypertarget\hypertarget{sourcepointer-2}{}\fi
\noindent\textit{Source reference.} Original lemma lemma-gerbe-residual-gerbe-exists: sources/stacks-algebraic-stacks-morphisms/pinned-original-source.tex, source lines 6684--6684. Navigation pointer only: the counted text refers to the Stacks item lemma-gerbe-residual-gerbe-exists, which is an existing labeled result of this same chapter and is not part of this card. The pointer carries the target's own label as its number and its definition line as its span; its content is not reproduced here.\par\medskip
\ifdefined\hypertarget\hypertarget{sourcepointer-3}{}\fi
\noindent\textit{Source reference.} Original lemma lemma-monomorphism-cartesian-square-inertia: sources/stacks-algebraic-stacks-morphisms/pinned-original-source.tex, source lines 1480--1480. Navigation pointer only: the counted text refers to the Stacks item lemma-monomorphism-cartesian-square-inertia, which is an existing labeled result of this same chapter and is not part of this card. The pointer carries the target's own label as its number and its definition line as its span; its content is not reproduced here.\par\medskip
\ifdefined\hypertarget\hypertarget{sourcepointer-4}{}\fi
\noindent\textit{Source reference.} Original proposition proposition-open-stratum: sources/stacks-algebraic-stacks-morphisms/pinned-original-source.tex, source lines 6746--6746. Navigation pointer only: the counted text refers to the Stacks item proposition-open-stratum, which is an existing labeled result of this same chapter and is not part of this card. The pointer carries the target's own label as its number and its definition line as its span; its content is not reproduced here.\par\medskip
\ifdefined\hypertarget\hypertarget{sourcepointer-5}{}\fi
\noindent\textit{Source reference.} Original reference stacks-properties-lemma-reduced-closed-substack: sources/stacks-algebraic-stacks-morphisms/pinned-original-source.tex, source lines 6854--6854. Navigation pointer only: the counted text refers to the Stacks item stacks-properties-lemma-reduced-closed-substack, which belongs to another chapter of the project and is not part of this card. The pointer carries the target's own label as its number because this chapter's staged snapshot does not contain the other chapter's tag table; its content is not reproduced here.\par\medskip
\begin{lemma}
\label{lemma-every-point-residual-gerbe}
Let $\mathcal{X}$ be an algebraic stack such that
$\mathcal{I}_\mathcal{X} \to \mathcal{X}$ is quasi-compact.
Then the residual gerbe of $\mathcal{X}$ at $x$ exists for
every $x \in |\mathcal{X}|$.
\end{lemma}

\begin{proof}
Let $T = \overline{\{x\}} \subset |\mathcal{X}|$ be the closure of $x$.
By
Properties of Stacks, Lemma
\ref{stacks-properties-lemma-reduced-closed-substack}
there exists a reduced closed substack $\mathcal{X}' \subset \mathcal{X}$
such that $T = |\mathcal{X}'|$. Note that
$\mathcal{I}_{\mathcal{X}'} =
\mathcal{I}_\mathcal{X} \times_\mathcal{X} \mathcal{X}'$ by
Lemma \ref{lemma-monomorphism-cartesian-square-inertia}.
Hence $\mathcal{I}_{\mathcal{X}'} \to \mathcal{X}'$ is
quasi-compact as a base change, see
Lemma \ref{lemma-base-change-quasi-compact}.
Therefore
Proposition \ref{proposition-open-stratum}
implies there exists a dense open substack
$\mathcal{U} \subset \mathcal{X}'$
which is a gerbe. Note that $x \in |\mathcal{U}|$ because $\{x\} \subset T$
is a dense subset too. Hence a residual gerbe
$\mathcal{Z}_x \subset \mathcal{U}$ of $\mathcal{U}$ at $x$ exists by
Lemma \ref{lemma-gerbe-residual-gerbe-exists}.
It is immediate from the definitions that $\mathcal{Z}_x \to \mathcal{X}$
is a residual gerbe of $\mathcal{X}$ at $x$.
\end{proof}

\end{document}
